\documentclass[conference]{IEEEtran}
\IEEEoverridecommandlockouts
\usepackage{cite}
\usepackage{amsmath,amssymb,amsfonts}
\usepackage{graphicx}
\usepackage{textcomp}
\usepackage{xcolor}
\usepackage{booktabs}   % clean tables (\toprule, \midrule, \bottomrule)
\usepackage{multirow}
\usepackage{url}

% If the paper needs Vietnamese text with diacritics (e.g. example tickets
% in the error-analysis table), test this line on Overleaf first:
% \usepackage[T5,T1]{fontenc}

% ---- Drafting helper: red notes. Delete all \TODO{...} before submission.
\newcommand{\TODO}[1]{\textcolor{red}{[TODO: #1]}}
\newcommand{\pmstd}[2]{#1{\scriptsize$\pm$#2}}

\begin{document}

% =====================================================================
% TITLE & AUTHORS
% =====================================================================
% Alternative title that matches the current data more closely:
% "Fine-Tuned Small Language Models for Robust Intent Classification of
%  Noisy Vietnamese Super-App Support Tickets"
\title{Fine-Tuning Small Language Models for Intent Classification of Vietnamese Customer Support Tickets}

\author{
\IEEEauthorblockN{\TODO{Member 1 full name}}
\IEEEauthorblockA{\textit{\TODO{Department}} \\
\textit{\TODO{University}}\\
\TODO{City}, Vietnam \\
\TODO{email}}
\and
\IEEEauthorblockN{\TODO{Member 2 full name}}
\IEEEauthorblockA{\textit{\TODO{Department}} \\
\textit{\TODO{University}}\\
\TODO{City}, Vietnam \\
\TODO{email}}
\and
\IEEEauthorblockN{\TODO{Member 3 full name}}
\IEEEauthorblockA{\textit{\TODO{Department}} \\
\textit{\TODO{University}}\\
\TODO{City}, Vietnam \\
\TODO{email}}
}

\maketitle

% =====================================================================
% ABSTRACT  (150-250 words, no citations / math / symbols)
% =====================================================================
\begin{abstract}
Vietnamese super-apps that combine ride-hailing, food and grocery delivery,
parcel delivery, e-wallets and membership services receive large volumes of
customer-support tickets that must be routed to the right team. Public
Vietnamese datasets for this task are scarce, and real tickets are often
written without diacritics, with chat abbreviations or with English terms.
We build a two-level intent taxonomy of 9 services and 26 issue types,
giving 45 valid labels, and a dataset of 900 synthetic tickets with six
test-time variants: clean text, diacritic removal, teencode, Telex typing
errors, code-switching and a mix of the three. We fine-tune PhoBERT-base and
Qwen2.5-Instruct models with 1.5 and 3 billion parameters using low-rank
adaptation, train only on clean text, and evaluate on every variant over
three random seeds. The 3-billion-parameter model reaches a macro-F1 of
90.1 percent on clean tickets against 81.0 percent for PhoBERT, and the gap
is significant on every variant. Diacritic removal is the most damaging
perturbation for all models, but PhoBERT keeps only about one fifth of its
clean-text score while the 3-billion-parameter model keeps nearly two
thirds. Tickets that do not name the service remain unsolved by all models.
The results suggest that small fine-tuned language models are a strong and
locally deployable option for routing noisy Vietnamese support tickets.
\TODO{Add one sentence on latency/cost once measured (Section~V-B).}
\end{abstract}

\begin{IEEEkeywords}
intent classification, small language models, LoRA, Vietnamese, customer support, noisy text, low-resource language
\end{IEEEkeywords}

% =====================================================================
\section{Introduction}
\label{sec:intro}
% =====================================================================

Vietnamese super-apps combine ride-hailing, food and grocery delivery,
parcel delivery, e-wallets and paid memberships in one application, so a
single support channel receives tickets about very different services.
Automatic intent classification routes each ticket to the responsible team
and shortens response time; customer service is one of the most common
deployment areas of AI \cite{fodor2026beyond}.

Three properties make this task hard in Vietnamese. First, tickets are
often non-standard: customers type without diacritics, use chat
abbreviations (teencode), leave Telex keystrokes in the text when the input
method is off, or mix in English terms such as \textit{refund} or
\textit{order}. Second, public Vietnamese intent datasets target chatbots
or flight booking rather than customer-support tickets
\cite{tran2020understanding}. Third, large API models are billed per call
and require sending tickets, which contain personal data, to an external
provider \TODO{cite Vietnam Personal Data Protection Law No.~91/2025/QH15 --
add to .bib}; small models that can be fine-tuned and served locally are
therefore attractive.

We study the following research questions:
\textbf{RQ1:} Do small generative language models fine-tuned with LoRA
outperform a fine-tuned Vietnamese encoder (PhoBERT) on super-app ticket
intents?
\textbf{RQ2:} How robust are these models to non-standard Vietnamese input
when they are trained on clean text only?
\textbf{RQ3:} What errors remain, and are they caused by the models or by
the data?
\TODO{Optional RQ4 on the accuracy--latency--cost trade-off, only if latency
is measured (Section~\ref{sec:results}-B).}

Our contributions are as follows:
\begin{itemize}
\item A two-level intent taxonomy for Vietnamese super-app support
(9 services $\times$ 26 issues, 45 valid labels) derived from public help
pages, and a 900-ticket synthetic dataset with six test-time variants of
non-standard Vietnamese.
\item A controlled comparison of PhoBERT-base (flat and hierarchical heads)
with Qwen2.5-1.5B and Qwen2.5-3B fine-tuned with LoRA and constrained
decoding, over three seeds and with paired bootstrap confidence intervals.
\item A robustness and error analysis showing that diacritic removal is the
dominant failure mode, that model scale mitigates it substantially, and
that tickets which do not name the service are not recoverable from the
text alone.
\end{itemize}

The rest of this paper is organized as follows. Section~\ref{sec:related}
reviews related work. Section~\ref{sec:dataset} describes dataset
construction. Section~\ref{sec:method} presents the models and experimental
setup. Section~\ref{sec:results} reports the results, and
Section~\ref{sec:discussion} discusses limitations. Section~\ref{sec:conclusion}
concludes the paper.

% =====================================================================
\section{Related Work}
\label{sec:related}
% =====================================================================
% Each sentence below is based on the team's Excel notes -- verify every
% claim against the original paper.

\subsection{Intent Detection for Vietnamese}
Tran and Luong \cite{tran2020understanding} built annotated Vietnamese
corpora for intent detection and context extraction in chatbots. Later work
studied joint intent detection and slot filling on PhoATIS, a Vietnamese
version of ATIS, under low-resource settings \cite{nguyen2022intent} and
with PhoBERT-based joint models \cite{nguyen2022advanced}. PhoBERT
\cite{nguyen2020phobert}, pre-trained on word-segmented Vietnamese text, is
the standard encoder baseline for such tasks. These studies target flight
booking or general chatbots; they do not address multi-service
customer-support tickets, small generative language models, or
non-standard input such as missing diacritics.

\subsection{Intent and Ticket Classification in Customer Service}
Machine-learning classifiers have been applied to customer-support
e-mails \cite{borg2021email}, banking chat requests
\cite{bonechi2024enhancing}, contact-centre conversations
\cite{pawlik2022method}, university administrative services
\cite{rizou2023efficient}, and cloud support tickets
\cite{nowakbrzezinska2025hybrid}. Other studies address discovering new
intents over time \cite{rebelo2022automatic}, refining intent sets with
linguistic analysis \cite{ferrera2025linguistics}, few-shot detection with
Siamese networks \cite{muhammad2025similarity}, explainability
\cite{licina2026helpagent}, multi-agent routing \cite{arslan2026application},
and how recognition errors affect users \cite{martijn2026answered}.
\TODO{One sentence on the common finding.} Most of this work concerns
English or European languages, and robustness to informal spelling is
rarely evaluated.

\subsection{Large and Small Language Models for Intent Classification}
LLMs have been used to generate training data for intent classifiers
\cite{benayas2024enhancing, madrueno2025novel} and compared with
weak supervision \cite{alexander2025few} and fine-tuned encoders
\cite{ouaddi2025assessing}. Fine-tuned small language models of 1.5B--4B
parameters have been evaluated for industrial customer-service tasks with
attention to efficiency \cite{pai2026finetuning}. Low-rank adaptation (LoRA)
\cite{hu2022lora} makes such fine-tuning feasible on a single consumer GPU,
and the Qwen2.5 family \cite{qwen2024qwen25} provides multilingual
instruction-tuned models at these sizes. To our knowledge, no study compares
LoRA-fine-tuned small language models with PhoBERT on Vietnamese
customer-support intents under controlled non-standard input.

% =====================================================================
\section{Dataset Construction}
\label{sec:dataset}
% =====================================================================

\subsection{Intent Taxonomy}
We first compiled a catalogue of 54 support intents in 9 categories from
the public help pages of Vietnamese super-apps; no internal company
documents were used. Many intents recur across services (e.g., being
charged more than the displayed price applies to rides, food, grocery and
parcel orders), so we factor each intent into a \emph{service} and an
\emph{issue}. The resulting two-level taxonomy has 9 services
(transport, food, mart, express, finance, business, membership, dine-out,
account) and 26 issue types; only 45 service--issue pairs are valid, and
the full label is the pair, e.g., \texttt{food/charged\_but\_failed}. Every
original intent maps to exactly one pair. For each issue, the annotation
guideline gives a definition, a ``not this'' rule that separates it from the
closest issue (e.g., \texttt{overcharged} vs.\ \texttt{charged\_but\_failed}),
the services it applies to, and example tickets. Two general rules apply:
a ticket with several problems receives the label of the main problem, and
a ticket that does not name the service is still assigned the most likely
service from context. The taxonomy and guideline are released with the
code \TODO{repository URL}.

\subsection{Ticket Generation}
Because public Vietnamese support tickets are not available, tickets were
written by an LLM (Claude \TODO{exact model version and date}) following the
guideline. The prompts asked for 20 tickets per label with varied forms of
self-address (\textit{t\^oi}, \textit{em}, \textit{m\`inh}), concrete details
(order codes, amounts, times, bank names) and both polite and frustrated
tones. Each ticket carries a flag \texttt{explicit} that is 0 when the
service is not stated (e.g., ``I was charged although the transaction
failed''). After whitespace normalization and duplicate removal (no
duplicates were found), the dataset contains 900 tickets, 20 per label.
\TODO{Manual check: who reviewed, how many tickets, rejection rate,
agreement (Cohen's kappa) -- currently not done; otherwise state it as a
limitation.}

\subsection{Splits}
We split the data 70/15/15, stratified by label, giving 14/3/3 tickets per
label for training, validation and test (Table~\ref{tab:data}). All
generation and splitting steps use a fixed seed (42).

\subsection{Noisy Variants}
Every ticket is rendered in six variants, generated by deterministic
rules: \textbf{clean} (original text); \textbf{no diacritics} (all tone and
vowel marks removed, and d-with-stroke replaced by d); \textbf{teencode}
(lower-casing and about twenty common chat abbreviations, e.g.,
\textit{kh\^ong}$\rightarrow$\textit{ko} ``not'',
\textit{m\`inh}$\rightarrow$\textit{mk} ``I'', and \textit{tk} for
``account''); \textbf{Telex typos}
(one or two words written as raw Telex keystrokes, as when the Vietnamese
input method is switched off); \textbf{code-switching} (support terms
replaced by their English equivalents such as \textit{refund},
\textit{cancel}, \textit{order}, \textit{payment}); and \textbf{mixed}
(code-switching, then teencode, then diacritic removal). Code-switching
changes only the 345 tickets (38\%) that contain a listed term; the other
variants change every ticket. Models are trained on clean text only and
evaluated on all six variants of the test set.

\subsection{Dataset Statistics}
Table~\ref{tab:data} summarizes the dataset. The number of tickets per
service ranges from 40 (membership, dine-out) to 160 (transport, food),
following the number of valid issues per service.

\begin{table}[htbp]
\caption{Dataset Statistics}
\label{tab:data}
\centering
\begin{tabular}{lrrrr}
\toprule
\textbf{Split} & \textbf{Tickets} & \textbf{Labels} & \textbf{Avg. words} & \textbf{Implicit} \\
\midrule
Train       & 630 & 45 & 12.2 & 23 \\
Validation  & 135 & 45 & 12.1 & 8 \\
Test        & 135 & 45 & 12.5 & 5 \\
\midrule
Total       & 900 & 45 & 12.2 & 36 \\
\bottomrule
\multicolumn{5}{p{0.92\columnwidth}}{\footnotesize Implicit: tickets that
do not name the service (\texttt{explicit}\,=\,0). The test set is
evaluated in six variants, i.e., $6\times135=810$ predictions per run.}
\end{tabular}
\end{table}

% =====================================================================
\section{Methodology}
\label{sec:method}
% =====================================================================

\subsection{Models}
\textbf{PhoBERT-base} (\texttt{vinai/phobert-base-v2})
\cite{nguyen2020phobert} is fully fine-tuned on word-segmented input
(segmentation with underthesea \cite{underthesea}). We use the first-token
representation with two head types: a \emph{flat} head over the 45 labels,
and a \emph{hierarchical} head with separate service (9) and issue (26)
outputs trained with the sum of two cross-entropy losses; at inference the
hierarchical model selects the valid pair $(s,i)$ that maximizes
$\log p(s)+\log p(i)$.

\textbf{Qwen2.5-1.5B-Instruct} and \textbf{Qwen2.5-3B-Instruct}
\cite{qwen2024qwen25} are fine-tuned with LoRA \cite{hu2022lora} as
instruction-to-label generators. The input is the model's chat template
with a short system prompt (``classify the support ticket of a Vietnamese
super-app; answer with one service/issue label'') and the ticket; the
target is the label string followed by the end-of-turn token. The loss is
computed on the label tokens only. At inference we use greedy decoding
constrained by a prefix tree over the tokenized 45 labels, so every output
is a valid label.

\TODO{Prompt-only LLM baseline (GPT-4o zero-/few-shot or an open model such
as Qwen2.5-7B-Instruct): add here if run. Note for a prompt-based baseline:
six example tickets in the annotation guideline appear verbatim in the test
set, so guideline examples must be removed from the prompt.}

\subsection{Training Configuration}
Table~\ref{tab:config} lists the hyper-parameters. Both model families use
an effective batch size of 16, linear decay with 10\% warm-up, gradient
clipping at 1.0, and model selection by macro-F1 on the clean validation
set after each epoch. Every configuration is run with seeds 42, 43 and 44.

\begin{table}[htbp]
\caption{Fine-Tuning Configuration}
\label{tab:config}
\centering
\setlength{\tabcolsep}{4pt}
\begin{tabular}{lll}
\toprule
 & \textbf{PhoBERT-base} & \textbf{Qwen2.5 + LoRA} \\
\midrule
Trainable params & 135M (all)  & 18.5M (1.5B), 29.9M (3B) \\
LoRA             & --          & $r=16$, $\alpha=32$, dropout 0.05 \\
LoRA targets     & --          & q, k, v, o, gate, up, down \\
Optimizer        & AdamW, wd 0.01 & AdamW, wd 0 \\
Learning rate    & $2\times10^{-5}$ & $2\times10^{-4}$ \\
Epochs           & 10          & 5 \\
Batch size       & 16          & 8 $\times$ 2 (grad.\ accum.) \\
Max length       & 64 tokens   & full prompt (no truncation) \\
Precision        & fp32        & bf16, no quantization \\
\bottomrule
\end{tabular}
\end{table}
% Trainable LoRA parameters computed from the Qwen2.5 configs
% (r=16 on all linear projections); check against
% model.print_trainable_parameters() in notebooks/C_qwen_lora.ipynb.

\subsection{Evaluation Metrics}
We report macro-F1 over the 45 full labels as the primary metric because it
weights all labels equally, together with accuracy and macro-F1 of the
service and issue parts. For generative models we also report the
invalid-output rate, i.e., the share of outputs that do not match any label;
with constrained decoding it is 0 for every run. Results are the mean
$\pm$ standard deviation over three seeds. Because the test set has only
three tickets per label, we compare each model with PhoBERT (hierarchical)
using a paired bootstrap over test tickets \cite{efron1993bootstrap}
(2{,}000 resamples, 95\% percentile intervals); for multi-seed models the
F1 on each resample is averaged over seeds.
\TODO{Latency (median and p95, ms per ticket), peak GPU memory, cost per
1{,}000 tickets -- not measured yet.}

\subsection{Implementation Details}
PhoBERT was trained on an Apple M4 Pro GPU (MPS backend) with PyTorch 2.8
and Transformers 4.57 (about 52\,s per run). The Qwen2.5 models were trained
on a single NVIDIA T4 GPU (Google Colab) with Transformers and PEFT
\TODO{versions; training time per run}. Decoding is greedy. Code, data and
predictions are available at \TODO{repository URL}.

% =====================================================================
\section{Results}
\label{sec:results}
% =====================================================================

\subsection{Main Results}
Table~\ref{tab:main} reports results on the clean test set. Both Qwen2.5
models outperform both PhoBERT variants on every metric. Qwen2.5-3B reaches
90.1\% macro-F1, 9.1 points above hierarchical PhoBERT, and Qwen2.5-1.5B
reaches 87.0\%. The bootstrap intervals of the gain over hierarchical
PhoBERT exclude zero for both Qwen models (Table~\ref{tab:boot}), whereas
the difference between the flat and hierarchical PhoBERT heads is not
significant ($-1.9$ points, CI $[-7.6, 3.2]$). The hierarchical head
improves service F1 for PhoBERT (91.5 vs.\ 88.9) but lowers issue F1
(84.4 vs.\ 87.6); the Qwen models improve both, with Qwen2.5-3B at 94.6 and
93.8.

\begin{table}[htbp]
\caption{Results on the Clean Test Set (\%, mean over 3 seeds)}
\label{tab:main}
\centering
\setlength{\tabcolsep}{3.5pt}
\begin{tabular}{lrcccc}
\toprule
\textbf{Model} & \textbf{Params} & \textbf{Acc.} & \textbf{Macro-F1} & \textbf{F1$_\text{svc}$} & \textbf{F1$_\text{issue}$} \\
\midrule
PhoBERT-base (flat)  & 135M & 80.0 & \pmstd{79.1}{1.1} & 88.9 & 87.6 \\
PhoBERT-base (hier.) & 135M & 82.2 & \pmstd{81.0}{1.4} & 91.5 & 84.4 \\
Qwen2.5-1.5B + LoRA  & 1.5B & 88.1 & \pmstd{87.0}{1.0} & 93.4 & 91.0 \\
Qwen2.5-3B + LoRA    & 3.1B & \textbf{90.6} & \textbf{\pmstd{90.1}{1.1}} & \textbf{94.6} & \textbf{93.8} \\
% GPT-4o zero-shot   & --   & -- & -- & -- & -- \\
% GPT-4o few-shot    & --   & -- & -- & -- & -- \\
\bottomrule
\end{tabular}
\end{table}

\subsection{Latency and Cost}
\TODO{Not measured yet. Measure on one GPU for all models: median/p95 ms per
ticket at batch size 1, peak memory, cost per 1{,}000 tickets. Remove this
subsection (and RQ4) if it will not be measured.}
% \begin{table}[htbp]
% \caption{Latency, Memory and Cost}
% \label{tab:cost}
% \centering
% \setlength{\tabcolsep}{4pt}
% \begin{tabular}{lccc}
% \toprule
% \textbf{Model} & \textbf{Latency (ms)} & \textbf{Memory (GB)} & \textbf{Cost / 1k tickets} \\
% \midrule
% PhoBERT-base         & -- & -- & -- \\
% Qwen2.5-1.5B + LoRA  & -- & -- & -- \\
% Qwen2.5-3B + LoRA    & -- & -- & -- \\
% \bottomrule
% \end{tabular}
% \end{table}

\subsection{Robustness to Non-Standard Input}
Table~\ref{tab:noise} and Fig.~\ref{fig:variants} show macro-F1 on all six
test variants. Teencode costs every model at most 2.4 points.
Code-switching costs PhoBERT 4.9--6.4 points but leaves both Qwen models
unchanged ($+0.7$ and $-0.1$), presumably because their pre-training covers
English. Telex typos cost PhoBERT 16--17 points but Qwen2.5-3B only 7.8. Diacritic removal
is the most damaging perturbation for all models: PhoBERT falls to 16\%
macro-F1, keeping about 20\% of its clean score, while Qwen2.5-1.5B keeps
40\% and Qwen2.5-3B keeps 63\% (56.3\%). The mixed variant, which includes
diacritic removal, behaves similarly (PhoBERT 17--19\%, Qwen2.5-3B 61.3\%).
Averaged over the five noisy variants, the drop from clean text is 31.0
points for hierarchical PhoBERT and 14.5 points for Qwen2.5-3B. All gains of
both Qwen models over PhoBERT are significant on every variant
(Table~\ref{tab:boot}). A plausible explanation is that PhoBERT relies on
word segmentation and on a vocabulary learned from standard, segmented text,
both of which break when diacritics are removed, whereas the byte-level
tokenizer and web-scale pre-training of Qwen2.5 expose it to unaccented
Vietnamese; the larger model benefits more.

\begin{table*}[t]
\caption{Macro-F1 (\%) on the Six Test Variants (mean $\pm$ std over 3 seeds). $\Delta$ = noisy average $-$ clean.}
\label{tab:noise}
\centering
\begin{tabular}{lcccccccc}
\toprule
\textbf{Model} & \textbf{Clean} & \textbf{No diac.} & \textbf{Teencode} & \textbf{Telex} & \textbf{Code-switch} & \textbf{Mixed} & \textbf{Noisy avg.} & $\boldsymbol{\Delta}$ \\
\midrule
PhoBERT-base (flat)  & \pmstd{79.1}{1.1} & \pmstd{16.0}{3.1} & \pmstd{78.2}{3.1} & \pmstd{63.2}{3.4} & \pmstd{74.2}{1.6} & \pmstd{19.3}{1.3} & \pmstd{50.2}{1.2} & $-28.9$ \\
PhoBERT-base (hier.) & \pmstd{81.0}{1.4} & \pmstd{16.1}{3.0} & \pmstd{78.6}{1.7} & \pmstd{64.2}{3.1} & \pmstd{74.6}{1.6} & \pmstd{16.6}{0.9} & \pmstd{50.0}{1.6} & $-31.0$ \\
Qwen2.5-1.5B + LoRA  & \pmstd{87.0}{1.0} & \pmstd{34.7}{5.6} & \pmstd{86.0}{0.8} & \pmstd{73.3}{1.5} & \pmstd{87.8}{0.7} & \pmstd{32.1}{2.2} & \pmstd{62.8}{1.4} & $-24.2$ \\
Qwen2.5-3B + LoRA    & \textbf{\pmstd{90.1}{1.1}} & \textbf{\pmstd{56.3}{1.4}} & \textbf{\pmstd{88.0}{2.1}} & \textbf{\pmstd{82.3}{0.7}} & \textbf{\pmstd{90.0}{1.2}} & \textbf{\pmstd{61.3}{1.2}} & \textbf{\pmstd{75.6}{0.4}} & $\mathbf{-14.5}$ \\
\bottomrule
\end{tabular}
\end{table*}

\begin{figure}[t]
\centering
\includegraphics[width=\columnwidth]{fig_f1_by_variant.pdf}
\caption{Macro-F1 on each test variant (training on clean text only).
Bars show $\pm$1 standard deviation over three seeds.}
\label{fig:variants}
\end{figure}

\begin{table}[htbp]
\caption{Macro-F1 Gain over PhoBERT-base (hier.) in Points, with 95\% Paired Bootstrap Intervals}
\label{tab:boot}
\centering
\setlength{\tabcolsep}{4pt}
\begin{tabular}{lcc}
\toprule
\textbf{Variant} & \textbf{Qwen2.5-1.5B + LoRA} & \textbf{Qwen2.5-3B + LoRA} \\
\midrule
Clean          & +6.1 [1.2, 13.0]   & +9.1 [4.6, 16.1] \\
No diacritics  & +18.6 [12.0, 22.4] & +40.3 [30.9, 44.1] \\
Teencode       & +7.5 [2.3, 13.8]   & +9.4 [4.3, 15.7] \\
Telex typos    & +9.1 [2.2, 17.3]   & +18.1 [12.4, 25.5] \\
Code-switching & +13.2 [8.2, 20.4]  & +15.4 [11.0, 22.7] \\
Mixed          & +15.5 [8.8, 19.7]  & +44.8 [35.8, 48.2] \\
\bottomrule
\end{tabular}
\end{table}
% Point estimate = observed difference of seed-mean macro-F1 (Table IV);
% interval = 2.5/97.5 percentiles of the paired bootstrap distribution
% (results/E_summary.md, Table 9, which lists bootstrap means instead).

\subsection{Error Analysis}
We classify each prediction as correct, wrong issue (service correct),
wrong service (issue correct), or both wrong (Fig.~\ref{fig:errors}). On
clean text the remaining errors of Qwen2.5-3B are few and spread across
types (2.5\% wrong issue, 4.4\% wrong service, 2.5\% both), whereas
hierarchical PhoBERT makes 7.7\% issue and 7.4\% service errors. Without
diacritics, 41\% of hierarchical PhoBERT predictions are wrong on both
levels, against 13\% for Qwen2.5-3B, i.e., PhoBERT loses the meaning of the
ticket rather than confusing neighbouring labels.

\begin{figure}[t]
\centering
\includegraphics[width=\columnwidth]{fig_error_types.pdf}
\caption{Error types on the clean and no-diacritics test variants, pooled
over three seeds.}
\label{fig:errors}
\end{figure}

The most frequent confusions on clean text are between food and mart
(missing or wrong items, undelivered and cancelled orders), payment issues
of other services predicted as \texttt{finance/charged\_but\_failed} by
PhoBERT, and \texttt{transport/trip\_howto} predicted as
\texttt{business/business\_profile} by both Qwen models. Six test tickets
are misclassified by all four systems; four of them are implicit tickets,
for example ``I was charged although the transaction reported failure''
(gold \texttt{mart/charged\_but\_failed}) and ``I cannot pay with my bank
card; no money was deducted'' (gold \texttt{mart/payment\_method\_failed}).
Nothing in these texts identifies the service, so the gold label cannot be
inferred. Accordingly, service accuracy on the five implicit test tickets
is 0\% for both PhoBERT variants, 6.7\% for Qwen2.5-1.5B and 13.3\% for
Qwen2.5-3B, against 89.7--96.2\% on explicit tickets. The two remaining
tickets contain misleading cues: ``Guide me to switch from my personal to
my business profile when booking a ride'' (gold
\texttt{transport/trip\_howto}) mentions the business profile, and ``Why
can't I see the driver moving on the map?'' (gold
\texttt{transport/trip\_howto}) is predicted as \texttt{driver\_no\_show}.
\TODO{If \texttt{[T5]\{fontenc\}} works on Overleaf, add the Vietnamese
originals of these examples.}

% =====================================================================
\section{Discussion and Limitations}
\label{sec:discussion}
% =====================================================================

\textbf{RQ1.} LoRA-fine-tuned Qwen2.5 models outperform PhoBERT on clean
super-app tickets, by 6.1 (1.5B) and 9.1 (3B) macro-F1 points over the
hierarchical head, while training only 18.5M--29.9M adapter parameters.
\textbf{RQ2.} When trained on clean text, all models are robust to
teencode and code-switching but not to diacritic removal; scale is the
main mitigating factor, since Qwen2.5-3B retains 63\% of its clean score
without diacritics against about 20\% for PhoBERT. For deployment, this
suggests either a model of at least 3B parameters or noise-aware training
\TODO{or diacritic restoration} when customers often type without
diacritics. \textbf{RQ3.} The residual errors on clean text are dominated by
tickets that do not name the service; such tickets should be resolved with
conversation context (e.g., the order the customer is viewing) or handed
to a human agent rather than by a larger classifier.

\textbf{Limitations.} (i) All tickets were written by a single LLM from
the annotation guideline, so their style is more uniform than real tickets,
and training and test data share the same generator; scores are likely
optimistic, and a human-written or real test set is needed. (ii) The tickets
have not yet been reviewed by human annotators \TODO{update if a review is
done}. (iii) Noisy variants are rule-based and may not match how customers
actually write; the test set and the noise rules come from the same
process. (iv) The test set has only three tickets per label and five
implicit tickets, so per-label scores are unreliable; we therefore report
only aggregate scores with bootstrap intervals. (v) Training budgets are not
fully matched: PhoBERT selected epoch 9--10 of 10 and Qwen2.5-1.5B epoch
4--5 of 5, so both may be under-trained, and the two families were trained
on different hardware. (vi) We consider single-label, single-turn tickets
only. (vii) Latency, memory and cost have not been measured yet, and no
prompt-only LLM baseline is included \TODO{update}.

% =====================================================================
\section{Conclusion}
\label{sec:conclusion}
% =====================================================================
We built a two-level intent taxonomy and a 900-ticket synthetic dataset for
Vietnamese super-app customer support with six variants of non-standard
input, and compared PhoBERT with Qwen2.5-1.5B and Qwen2.5-3B fine-tuned with
LoRA. Qwen2.5-3B achieved 90.1\% macro-F1 on clean tickets and 75.6\% on
average over noisy variants, against 81.0\% and 50.0\% for PhoBERT, with the
largest gains on tickets without diacritics. Future work includes
evaluating on real or human-written tickets, noise-aware training and
diacritic restoration, measuring latency for deployment, multi-label
intents, confidence-based hand-off to human agents, and multi-turn context
for tickets that do not name the service.

% =====================================================================
\section*{Acknowledgment}
% =====================================================================
\TODO{Thank the supervisor. Delete this section if not needed.}

% =====================================================================
% REFERENCES -- uses references.bib (upload it to the Overleaf project)
% =====================================================================
\bibliographystyle{IEEEtran}
\bibliography{references}

\end{document}
